% Part: lambda-calculus
% Chapter: church-rosser
% Section: definitions-and-properties

\documentclass[../../../include/open-logic-section]{subfiles}

\begin{document}

\olfileid{lam}{cr}{dap}

\olsection{تعریف تے خاصیتاں}

ایس باب وچ اسیں چرچ–روسر خاصیت دا خیال تے ایہدیاں کجھ
عام خاصیتاں پیش کردے آں۔

\begin{defn}[چرچ–روسر خاصیت، CR]
  ٹرماں اُتے تعلق $\xredone$ نوں \emph{چرچ–روسر خاصیت} والا
  اُدوں تے صرف اُدوں آکھدے آں جدوں $M \xredone P$ تے
  $M \xredone Q$ ہون دی ہر صورت وچ کوئی~$N$ موجود ہووے
  جیہدے لئی $P \xredone N$ تے $Q \xredone N$۔
\end{defn}

اسیں لامبڈا کیلکولس نوں حساب دا اک نمونہ سمجھ سکدے آں،
جیہدے وچ نارمل روپ والے ٹرم «قدراں» نیں تے ٹرماں اُتے
تقلیل دا تعلق «حساب دے قاعدے» دیندا اے۔ چرچ–روسر خاصیت
آکھدی اے کہ حساب اگّے ودھاؤن دے کئی طریقے ہون تے وی عبارت
دی اکّو قدر ہوندی اے۔

مُڈھلے الجبرا دی اک مثال ویکھو۔ $4 \times (1+2) + 3$ دا
حساب اک توں ودھ طریقیاں نال ہو سکدا اے۔ ایہدی تقلیل
$4
\times 3+3$ تک ہو سکدی اے، جے پہلاں $1+2$ دی تقلیل
~$3$ تک کریے۔ یا $4 \times 1+4
\times 2+3$ تک ہو سکدی اے، جے پہلاں $4 \times (1+2)$ دی
تقلیل ضرب نوں جمع دے دوہاں حصیاں اُتے ونڈن دے قاعدے نال کریے۔ دوہاں دی اگلی تقلیل
$12+3$ تک ہو سکدی اے۔

جے $\xredone$ دی تھاں $\beta$-تقلیل لئیے، تاں سوکھا ای
دِسدا اے کہ چرچ–روسر خاصیت توں ایہ نتیجہ نکلدا اے کہ جے
کسے ٹرم دا نارمل روپ موجود ہووے، تاں اوہ اکّو اے۔ فرض کرو
$M$ دی تقلیل $P$ تے $Q$ تک ہو سکدی اے، تے دوویں نارمل
روپ نیں۔ چرچ–روسر خاصیت مطابق کوئی $N$ موجود اے جیہدے
تک $P$ تے $Q$ دوہاں دی تقلیل ہوندی اے۔ مفروضے مطابق $P$
تے $Q$ نارمل روپ نیں، ایس لئی $P$ تے $Q$ دی $N$ تک
تقلیل صرف صفر قدماں والی تقلیل ہو سکدی اے۔ یعنی $P$،
$Q$ تے $N$ اکّو ٹرم نیں۔ ایس کرکے ٹرم دے \emph{اکّو}
نارمل روپ دی گل کرنا ٹھیک اے۔

سو جدوں اسیں لامبڈا کیلکولس نوں حساب دا نمونہ سمجھدے آں،
تاں کسے ٹرم دا نارمل روپ اوہ حساب دا «آخری نتیجہ» سمجھیا
جا سکدا اے جیہڑا ایس ٹرم توں شروع ہوندا اے۔ اُپرلا نتیجہ
دسدا اے کہ جے حساب دا کوئی آخری نتیجہ ہووے، تاں اکّو
ہووے گا۔ جیویں $4 \times (1+2)+3$ دا نتیجہ اکّو اے،
یعنی~$15$۔

\begin{thm} \ollabel{thm:str}
  جے تعلق $\xredone$ وچ چرچ–روسر خاصیت اے، تے $\xred$
  سبھ توں چھوٹا تعدّی والا تعلق اے جیہدے وچ $\xredone$
  شامل اے، تاں $\xred$ وچ وی چرچ–روسر خاصیت اے۔
\end{thm}

\begin{proof}
  فرض کرو ہِٹھلیاں زنجیراں نیں۔
  \begin{align*}
    M & \xredone P_1 \xredone \dots \xredone P_m \text{ تے}\\
    M & \xredone Q_1 \xredone \dots \xredone Q_n.
  \end{align*}
  قضیہ ثابت کرن لئی اسیں ٹرماں دا جال $N$ بناواں گے،
  جیہدی اُچائی $m + 1$ تے چوڑائی $n + 1$ اے۔ $N_{i,j}$
  نال $i$ویں قطار تے $j$ویں کالم دا ٹرم لکھدے آں۔

  اسیں $N$ ایویں بناواں گے کہ $N_{i,j} \xredone N_{i+1,j}$ تے
  $N_{i,j} \xredone N_{i,j+1}$۔ تعریف ایویں اے۔
  \begin{align*}
    N_{0,0} &= M \\
    N_{i,0} &= P_i && \text{جے } 1 \le i \le m \\
    N_{0,j} &= Q_j && \text{جے } 1 \le j \le n \\
  \intertext{تے ہور صورت وچ}
    N_{i,j} &= R
  \end{align*}
  ایتھے $R$ اوہ ٹرم اے جیہدے لئی $N_{i-1,j} \xredone R$ تے
  $N_{i,j-1}
  \xredone R$۔ $\xredone$ دی چرچ–روسر خاصیت مطابق ایہو
  جیہا ٹرم ہمیشہ موجود ہوندا اے۔

  ہُن ساڈے کول $N_{m,0} \xredone \dots \xredone N_{m,n}$ تے
  $N_{0,n}
  \xredone \dots \xredone N_{m,n}$ نیں۔ دھیان دیو کہ $N_{m,0}$،
  $P$ اے تے $N_{0,n}$، $Q$ اے۔ $\xred$ دی تعریف توں
  قضیہ نکلدا اے۔
\end{proof}

\end{document}
